% ============================================================
%  Chapter 6 — Conclusion and Future Work
% ============================================================
\chapter{Conclusion and Future Work}
\label{chap:conclusion}

This thesis proposed DRL-LSTM-LBRO, a framework that augments
the heuristic Load Balancing with Resource Optimisation (LBRO)
strategy with proactive GRU-LSTM load forecasting and a
Double Deep Q-Network (DDQN) scheduling agent for adaptive
task placement in IoT edge computing environments.

% ─────────────────────────────────────────────────────────
\section{Summary of Contributions}
\label{sec:summary}

The principal contributions of this work are as follows:

\begin{enumerate}

  \item \textbf{GRU-LSTM Load Predictor.}
  A per-cloudlet GRU-LSTM model is trained to forecast
  near-future CPU and RAM utilisation from a sliding window
  of $L = 10$ historical observations. The two-stage
  recurrent architecture — a GRU layer for short-term
  dynamics followed by an LSTM layer for longer-range
  patterns — produces accurate 1-step utilisation forecasts
  that are incorporated into the DDQN state vector, enabling
  the agent to act on future load trends rather than only
  current measurements.

  \item \textbf{Multi-Objective DDQN Scheduling Agent.}
  A DDQN agent with separate online and target networks is
  designed to select the optimal task-placement action from
  a discrete action space covering $N$ heterogeneous
  cloudlets and a cloud fallback. The reward function jointly
  penalises task drops, latency, energy consumption, and load
  imbalance, enabling the agent to balance competing
  objectives without manual tuning of a fixed CRI formula.

  \item \textbf{End-to-End Integration with LBRO.}
  The framework extends the original LBRO heuristic by
  replacing its static Composite Resource Index with the
  learned DDQN policy. By retaining the LBRO-inspired
  problem structure (finite-capacity cloudlets, M/M/$c$
  queuing model, cloud fallback) while replacing the
  decision rule with a data-driven agent, the framework
  achieves strong performance without discarding the
  domain knowledge embedded in LBRO.

  \item \textbf{Empirical Evaluation on Real Workload Traces.}
  Experiments on the Google Borg cluster trace (294,725
  tasks) with 10 random seeds and 50 evaluation episodes
  per seed demonstrate that DRL-LSTM-LBRO reduces task drop
  rate by 79\% relative to LBRO (0.58\% vs.\ 2.76\%) and
  average latency by 23.4\% (0.2603\,s vs.\ 0.3397\,s).
  Ablation confirms an 8.5\% additional latency reduction
  attributable to GRU-LSTM state augmentation. Scalability
  tests on a 6-cloudlet topology show the framework
  generalises without architectural modification.

\end{enumerate}

% ─────────────────────────────────────────────────────────
\section{Limitations}
\label{sec:limitations}

Despite the encouraging results, the current work has several
limitations that should be acknowledged:

\begin{enumerate}

  \item \textbf{Load Imbalance Trade-off.}
  The DDQN policy learns to concentrate tasks on the
  fastest cloudlets to maximise the latency and drop-rate
  reward components. This results in a lower Jain's Fairness
  Index (0.8762 on 3 cloudlets; 0.6759 on 6 cloudlets)
  compared to heuristic baselines ($>$0.99). While the
  agent's behaviour is consistent with the reward design,
  deployments requiring strict fairness guarantees would
  need a higher imbalance penalty coefficient or a
  constrained RL formulation.

  \item \textbf{Simulation-Only Evaluation.}
  All experiments are conducted in a discrete-event
  simulator. Although the workload traces originate from
  a real Google data-centre cluster, the simulator
  abstracts away network heterogeneity, hardware failures,
  and operating-system scheduling variability. Real-world
  deployment may exhibit additional noise sources.

  \item \textbf{Static Topology.}
  The current framework assumes a fixed set of cloudlets
  known at training time. Node joins, failures, or
  dynamic topology changes require the DDQN to be
  retrained or fine-tuned, which may be impractical in
  highly volatile edge environments.

  \item \textbf{Single-Step Forecasting.}
  The GRU-LSTM predictor provides a 1-step-ahead
  utilisation forecast. Multi-step prediction could
  allow the agent to plan further ahead during burst
  arrivals, but increases model complexity and
  potential compounding forecast error.

\end{enumerate}

% ─────────────────────────────────────────────────────────
\section{Future Work}
\label{sec:future}

Building on the contributions and limitations above, the
following directions are proposed for future research:

\begin{enumerate}

  \item \textbf{Constrained RL for Fairness.}
  Incorporating a hard fairness constraint into the RL
  optimisation (e.g., via Lagrangian relaxation or
  constrained policy optimisation) would allow the agent
  to achieve low drop rate and latency while guaranteeing
  a minimum JFI, making the framework suitable for
  service-level agreement (SLA)-bound deployments.

  \item \textbf{Multi-Step and Probabilistic Forecasting.}
  Extending the GRU-LSTM to produce $k$-step-ahead
  probabilistic forecasts (e.g., with Monte Carlo dropout
  or conformal prediction intervals) would give the agent
  uncertainty-aware state information, enabling more robust
  decisions during high-variance burst traffic periods.

  \item \textbf{Dynamic Topology Adaptation.}
  Applying transfer learning or meta-RL techniques to
  allow the DDQN agent to adapt quickly to topology
  changes (new cloudlets, node failures) without
  full retraining would be a significant step toward
  production-ready deployment.

  \item \textbf{Real Hardware Testbed Deployment.}
  Validating DRL-LSTM-LBRO on a physical edge testbed
  (e.g., Raspberry Pi cluster or OpenStack-based cloudlet
  deployment) would measure the framework's performance
  under realistic networking delays, CPU contention, and
  OS scheduling jitter.

  \item \textbf{Federated and Privacy-Preserving Learning.}
  In multi-operator edge environments, sharing raw
  workload traces for centralised LSTM training may
  violate data-privacy agreements. Federated learning
  approaches that train the GRU-LSTM locally at each
  cloudlet and aggregate model updates without exchanging
  raw data represent a practical path to wider deployment.

  \item \textbf{Hierarchical Scheduling.}
  Combining a coarse-grained global DDQN agent (selecting
  which cloudlet cluster to use) with a fine-grained
  local agent (scheduling tasks within a cloudlet) could
  reduce the action space at each level, accelerating
  convergence on large-scale topologies with tens of nodes.

\end{enumerate}

In summary, DRL-LSTM-LBRO demonstrates that integrating
predictive load awareness with end-to-end reinforcement
learning is a viable and effective approach to adaptive
load balancing in IoT edge computing. The 79\% drop rate
reduction and 23\% latency improvement over the LBRO
baseline confirm the value of replacing hand-crafted
heuristics with learned, forecast-augmented policies, and
establish a clear foundation for the future research
directions outlined above.
